\section*{Bab \ref{ch:b2:heat-conduction} --- Hantaran Kalor}

\begin{solution}{exo:b2:heat-conduction:1}
$j = \lambda\Delta T/e = \qty{80}{W/m^2}$; \qty{8}{kW}; dengan wolnya $R = 0.25 + 2.5 =
\qty{2.75}{m^2.K/W}$: \qty{7.3}{W/m^2}, \qty{730}{W}.
\end{solution}

\begin{solution}{exo:b2:heat-conduction:2}
$a$: \qty{1.2e-4}{m^2/s}, \qty{7e-7}{m^2/s}, \qty{1e-7}{m^2/s}. \qty{1}{cm}:
\qty{0.9}{s}, \qty{140}{s}, \qty{1000}{s}; \qty{20}{cm}: \qty{6}{min}, \qty{15}{h},
\qty{5}{hari}; keraknya ($a \sim \qty{1e-6}{m^2/s}$): $10^{15}\,\unit{s}$, tiga puluh
juta tahun.
\end{solution}

\begin{solution}{exo:b2:heat-conduction:3}
Tunggal: $R = 0.169$, $U = \qty{5.9}{W/m^2/K}$, \qty{118}{W/m^2}; ganda: $R =
0.63$, $U = \qty{1.6}{W/m^2/K}$, \qty{32}{W/m^2}. Celahnya berkonveksi dan kacanya
bertukar lewat radiasi (sekitar separuh alih yang sebenarnya); argon menghantar
lebih sedikit dan berkonveksi lebih sedikit, sebuah lapisan beremisivitas rendah menekan
radiasinya.
\end{solution}

\begin{solution}{exo:b2:heat-conduction:4}
$L^2/a$ dengan setengah tebalnya: \qty{27}{min}; panggangnya: $16\times$, \qty{7}{h};
$t \propto L^2 \propto m^{2/3}$: $\times 1.6$.
\end{solution}

\begin{solution}{exo:b2:heat-conduction:5}
(a) $(rT')'/r = -p/\lambda$ dengan $T'(0) = 0$. (b) $p = j^2/\gamma = \qty{1.7e6}{W/m^3}$;
$pR^2/4\lambda = \qty{1e-3}{K}$. (c) \qty{625}{K}. (d) Sebuah kawat itu isotermal; pusat
sebuah pelet bahan bakar beratus kelvin di atas permukaannya --- rapat daya
sebuah reaktor dibatasi hantaran dalam bahan bakarnya.
\end{solution}

\begin{solution}{exo:b2:heat-conduction:6}
(a) \cref{prop:b2:heat-conduction:wave}. (b) \qty{12}{cm}; \qty{2.2}{m}. (c)
$\delta\ln 10 = \qty{5.2}{m}$; tundanya $2.3\,\unit{rad}$, $130$ hari. (d) Amplitudo
\qty{20}{K}-nya harus menjadi separuh: $x = \delta\ln 2 = \qty{1.5}{m}$.
\end{solution}

\begin{solution}{exo:b2:heat-conduction:7}
(a) $\mathrm{Bi} = hR/3\lambda \approx 10^{-4}$; $\tau = \rho cR/3h = \qty{1100}{s}$.
(b) $\mathrm{Bi} = 1.2$: bukan salah satu limitnya; waktu dalamnya $R^2/a \approx
\qty{6000}{s}$ yang mengatur. (c) $R$ kecil: $\tau$ mungil dan \omterm{def:b2:heat-conduction:biot}{bilangan Biot} kecil.
(d) Hambatan selaputnya $1/h$, yang menguasai, turun.
\end{solution}

\begin{solution}{exo:b2:heat-conduction:8}
(a) \cref{prop:b2:heat-conduction:fin}. (b) $m = \sqrt{2h/\lambda t} =
\qty{13}{m^{-1}}$, $mL = 0.39$, $\eta = 0.95$. (c) Luasnya $\times 2L/t = 60$, fluksnya
$\times 57$. (d) Di luar $mL \approx 2$ panjang tambahannya berada pada suhu
fluidanya; sirip yang lebih tipis mempunyai $m$ yang lebih besar dan efisiensi yang lebih rendah.
\end{solution}

\begin{solution}{exo:b2:heat-conduction:9}
(a) \cref{prop:b2:heat-conduction:resistance}. (b) Wolnya $\ln 2/2\pi\lambda =
\qty{2.76}{K.m/W}$, selaput luarnya $1/2\pi r_2h = 0.16$: \qty{45}{W/m}; telanjang,
$2\pi r_1h\Delta T = \qty{410}{W/m}$. (c) $\dd/\dd r[\ln(r/r_1)/2\pi\lambda + 1/2\pi hr]
= 0$ pada $r = \lambda/h$. (d) $r_{\text{c}} = \qty{2}{cm} > \qty{1}{mm}$: selubungnya
menaikkan ruginya --- ia mendinginkan tembaganya.
\end{solution}

\begin{solution}{exo:b2:heat-conduction:10}
(a) $8000(1/273 - 1/293) = \qty{2.0}{W/K}$. (b) $\sigma_s = \lambda T'^2/T^2$ dengan
$T' = \qty{100}{K/m}$; $\int\sigma_s\,\dd x = \lambda T'(1/T_{\text{c}} - 1/T_{\text{h}})
= j(1/T_{\text{c}} - 1/T_{\text{h}})$; dikalikan \qty{100}{m^2}: sama saja. (c)
\qty{0.18}{W/K}. (d) $\Phi(1 - T_{\text{c}}/T_{\text{h}}) = \qty{550}{W}$; tak ada apa pun
--- kalornya kini duduk pada $T_{\text{c}}$ ($T_{\text{c}}\dot S_{\text{c}} = \qty{550}{W}$).
\end{solution}

\begin{solution}{exo:b2:heat-conduction:11}
(a) $j_Q(0,t) = -\lambda\partial_xT = \lambda\theta_0\sqrt2/\delta\,\cos(\omega t +
\pi/4) = \sqrt{\lambda\rho c\omega}\,\theta_0\cos(\omega t + \pi/4)$. (b) Fluks tiap
sisinya di sentuhannya adalah $b_i|T_i - T_{\text{c}}|/\sqrt{\pi t}$; samakanlah:
$T_{\text{c}} = (b_1T_1 + b_2T_2)/(b_1 + b_2)$. (c) Tembaga \qty{20.3}{\celsius},
kayu \qty{29}{\celsius}. (d) Ubin $b \approx 1500$, karpet $\approx 100$;
rumusnya berlaku selama kedua bendanya tampak setengah takhingga; begitu muka
kalornya mencapai pasokan darahnya dan sisi jauh bendanya, keadaan
tunaknya bergantung pada kalor yang sungguh dipasok.
\end{solution}

\begin{solution}{exo:b2:heat-conduction:12}
(a) Sisipkanlah, lalu $\theta(\pm L) = 0$ memberi $\cos k_nL = 0$. (b) $a_n =
4\theta_0(-1)^n/(2n+1)\pi$; $\theta = \sum a_n\cos(k_nx)\eu^{-ak_n^2t}$. (c) $n = 0$,
$\tau = 4L^2/\pi^2a$. (d) $\tau = \qty{650}{s}$; pusatnya $(4/\pi)\eu^{-t/\tau} = 0.01$:
$t = 4.8\tau \approx \qty{52}{min}$.
\end{solution}

\begin{solution}{pb:b2:heat-conduction:1}
\textbf{1.} \qty{0.25}{m^2.K/W}; \qty{80}{W/m^2}.

\textbf{2.} $R = 0.415$, $U = \qty{2.4}{W/m^2/K}$; yang tersekat $R = 2.9$, $U =
\qty{0.34}{W/m^2/K}$.

\textbf{3.} \qty{5.9}{W/m^2/K}; \qty{1.6}{W/m^2/K}.

\textbf{4.} $4800 + 2360 = \qty{7.2}{kW}$; $680 + 640 = \qty{1.3}{kW}$.

\textbf{5.} $UA = \qty{358}{W/K}$: $358 \times 2500 \times 86400 = \qty{7.7e10}{J} =
\qty{21500}{kWh}$; sesudahnya, \qty{66}{W/K}: \qty{4000}{kWh}.

\textbf{6.} $20 - 48/8 = \qty{14}{\celsius}$; $20 - 6.8/8 = \qty{19.2}{\celsius}$;
dinding yang dingin memancar lebih sedikit ke tubuhnya dan dapat mencapai titik embunnya.

\textbf{7.} $a = \qty{5.3e-7}{m^2/s}$; $e^2/a = \qty{7.5e4}{s} \approx \qty{21}{h}$:
dindingnya merata-ratakan harinya (daur harian $\delta = \qty{12}{cm}$ lebih kecil daripada
tebalnya).

\textbf{8.} Carilah $\theta_0\eu^{\iu(\omega t - \underline kx)}$: $\underline k = (1 - \iu)/\delta$.

\textbf{9.} \qty{12}{cm}; \qty{2.2}{m}.

\textbf{10.} $4.6\delta = \qty{54}{cm}$; \qty{3.7}{K}; \qty{1.1}{K}.

\textbf{11.} Satu radian, $58$ hari: pertengahan Maret (karena minimum
permukaannya pertengahan Januari).

\textbf{12.} Amplitudonya $\lambda\theta_0\sqrt2/\delta = \qty{6.3}{W/m^2}$, yang mendahului sebesar
$\pi/4$ ($45$ hari); seratus kali fluks geotermalnya.

\textbf{13.} $\lambda \times \qty{0.03}{K/m} \approx \qty{0.03}{W/m^2}$, ordenya benar; ia
merupakan \qty{0.1}{K} yang tunak sepanjang tinggi ruang bawah tanahnya, terkubur dalam
ayunan musimannya.

\textbf{14.} \qty{1e6}{W/m^2}; $R = \qty{0.033}{K/W}$.

\textbf{15.} \qty{0.01}{K/W}; udaranya: \qty{1.9}{K/W}, \qty{190}{K}.

\textbf{16.} $m = \qty{20}{m^{-1}}$, $mL = 0.61$, $\eta = 0.89$; $\mathrm{Bi} = ht/2\lambda
= 10^{-4}$.

\textbf{17.} $A = \qty{0.108}{m^2}$; $R_{\text{sirip}} = 1/\eta hA = \qty{0.21}{K/W}$; alasnya
\qty{0.006}{K/W}; totalnya \qty{0.26}{K/W}: \qty{26}{K}, sambungannya pada \qty{51}{\celsius}.

\textbf{18.} $1/hS = \qty{5.6}{K/W}$: \qty{560}{K} --- mustahil.

\textbf{19.} $R_{\text{sirip}} \approx \qty{0.94}{K/W}$, totalnya $\approx \qty{1}{K/W}$:
\qty{125}{\celsius} --- ia menurunkan lajunya, lalu mati.

\textbf{20.} $V = \qty{7.2e-5}{m^3}$, \qty{0.19}{kg}, $C = \qty{175}{J/K}$; $\tau = RC
\approx \qty{46}{s}$; kepingnya: $C = \qty{0.07}{J/K}$, $\tau \approx \qty{3}{ms}$.

\textbf{21.} Kalornya harus menyebar dari \qty{1}{cm^2} ke \qty{36}{cm^2}: tembaga
memisahduakan hambatan penyebaran itu; sebuah pipa kalor memindahkan kalor lewat
penguapan dan pengembunan dengan konduktivitas efektif seratus
kali milik tembaga.

\textbf{22.} $7200(1/273 - 1/293) = \qty{1.8}{W/K}$; \qty{0.33}{W/K}.

\textbf{23.} $100(1/298 - 1/324) = \qty{0.027}{W/K}$; $100/324 = \qty{0.31}{W/K}$ ---
sepuluh kali lebih banyak: mengubah kerja menjadi kalor adalah ketakterbalikan yang besar.

\textbf{24.} $\int\lambda T'^2/T^2\,\dd x = \lambda T'\,[-1/T] = j(1/T_{\text{c}} -
1/T_{\text{h}})$.

\textbf{25.} Dinding: hambatan yang terseri, penghantar terburuknya yang berkuasa.
Sirip: $m = \sqrt{hp/\lambda S}$, $mL \approx 1$. Ruang bawah tanah: $\delta = \sqrt{2a/\omega}$,
peredaman dan tunda bersama-sama. Cip: sebuah rantai hambatan dari kepingnya ke
udaranya, dengan luas siripnya yang bekerja. Di mana-mana: kalor mengalir menurun, lalu
menciptakan entropi sambil melakukannya.
\end{solution}
